\chapter{Programs — Regions}
%%% generated by the blueprint pipeline from TCSlib.Programs.Light.Regions
\section{Overview}

Declarations of the Lean module \texttt{TCSlib.Programs.Light.Regions}.

\section{Declarations}

\begin{definition}[Membership in a natural-number interval]
\lean{Light.Inside}
\label{Light.Inside}
\leanok
For natural numbers $a,n,b$, say that $b$ is inside the interval starting at $a$ with
length $n$ when $a\le b<a+n$.
\end{definition}

\begin{definition}[Outside a natural-number interval]
\lean{Light.Outside}
\label{Light.Outside}
\leanok
For natural numbers $a,n,b$, say that $b$ is outside the interval starting at $a$ with
length $n$ when $b<a$ or $a+n\le b$.
\end{definition}

\begin{definition}[Disjoint ordered intervals]
\lean{Light.Apart}
\label{Light.Apart}
\leanok
For natural numbers $a,n,a',n'$, the intervals starting at $a$ and $a'$ with lengths $n$
and $n'$ are apart when $a+n\le a'$ or $a'+n'\le a$.
\end{definition}

\begin{definition}[Ordered list of intervals]
\lean{Light.InOrder}
\label{Light.InOrder}
\leanok
A list of natural-number pairs is in order below $T$ if it is empty, if its sole pair
$(a,n)$ satisfies $a+n\le T$, or, when it has at least two pairs, each pair's starting
point plus length is at most the next pair's starting point and this condition holds
recursively for the remaining list.
\end{definition}

\begin{definition}[Agreement on a set of indices]
\lean{Light.SameOn}
\label{Light.SameOn}
\leanok
For a predicate $K$ on the natural numbers and integer-valued sequences $\mu,\mu'$, say
that they agree on $K$ when $\mu'(b)=\mu(b)$ for every natural number $b$ satisfying
$K(b)$.
\end{definition}

\begin{definition}[Agreement outside an interval]
\lean{Light.SameOutside}
\label{Light.SameOutside}
\leanok
\uses{Light.Outside, Light.SameOn}
For natural numbers $a,n$ and integer-valued sequences $\mu,\mu'$, say that they agree
outside the interval specified by $a,n$ when they agree at every index outside that
interval.
\end{definition}

\begin{definition}[Agreement outside two intervals]
\lean{Light.SameOutside2}
\label{Light.SameOutside2}
\leanok
\uses{Light.Outside, Light.SameOn}
For natural numbers $a,n,b,m$ and integer-valued sequences $\mu,\mu'$, say that they
agree at every index that lies outside both intervals specified by $a,n$ and $b,m$.
\end{definition}

\begin{definition}[Agreement outside three intervals]
\lean{Light.SameOutside3}
\label{Light.SameOutside3}
\leanok
\uses{Light.Outside, Light.SameOn}
For natural numbers $a,n,b,m,c,k$ and integer-valued sequences $\mu,\mu'$, say that they
agree at every index that lies outside all three intervals specified by $a,n$, $b,m$,
and $c,k$.
\end{definition}

\begin{definition}[Agreement below a cutoff]
\lean{Light.Kept}
\label{Light.Kept}
\leanok
\uses{Light.SameOn}
For integer-valued sequences $\mu,\mu'$ and a natural number $fr$, say that they agree
below $fr$ when $\mu'(x)=\mu(x)$ for every natural number $x<fr$.
\end{definition}

\begin{definition}[Agreement below a cutoff outside an interval]
\lean{Light.KeptBut}
\label{Light.KeptBut}
\leanok
\uses{Light.Outside, Light.SameOn}
For integer-valued sequences $\mu,\mu'$ and natural numbers $fr,out,len$, say that they
agree at every index $x<fr$ that lies outside the interval specified by $out,len$.
\end{definition}

\begin{theorem}[Agreement at a specified cell]
\lean{Light.SameOn.cell}
\label{Light.SameOn.cell}
\leanok
\uses{Light.SameOn}
\difficulty{1}
Let $K:\mathbb{N}\to\mathrm{Prop}$ be a predicate and let
$\mu,\mu':\mathbb{N}\to\mathbb{Z}$ agree at every index satisfying $K$. If $\mu$ and
$\mu'$ agree at exactly the index $b$, then $\mu'(b)=\mu(b)$.
\end{theorem}

\begin{theorem}[Reflexivity of agreement]
\lean{Light.SameOn.refl}
\label{Light.SameOn.refl}
\leanok
\uses{Light.SameOn}
\difficulty{1}
For any predicate $K$ on the natural numbers and any integer-valued sequence $\mu$, the
sequences $\mu$ and $\mu$ agree at every index satisfying $K$.
\end{theorem}

\begin{theorem}[Transitivity of agreement]
\lean{Light.SameOn.trans}
\label{Light.SameOn.trans}
\leanok
\uses{Light.SameOn}
\difficulty{1}
Let $K$ be a predicate on the natural numbers, and let $\mu,\mu',\mu''$ be
integer-valued sequences. If $\mu'$ agrees with $\mu$ and $\mu''$ agrees with $\mu'$ at
every index satisfying $K$, then $\mu''$ agrees with $\mu$ at every such index.
\end{theorem}

\begin{theorem}[Monotonicity of memory agreement]
\lean{Light.SameOn.mono}
\label{Light.SameOn.mono}
\leanok
\uses{Light.SameOn}
\difficulty{1}
Let $K,K':\mathbb{N}\to\mathrm{Prop}$ and let $\mu,\mu':\mathbb{N}\to\mathbb{Z}$. If
$\mu$ and $\mu'$ agree at every index satisfying $K$, and every index satisfying $K'$
also satisfies $K$, then $\mu$ and $\mu'$ agree at every index satisfying $K'$.
\end{theorem}

\begin{theorem}[Composition of agreement]
\lean{Light.SameOn.then}
\label{Light.SameOn.then}
\leanok
\uses{Light.SameOn}
\difficulty{1}
Let $K,K_1,K_2$ be predicates on the natural numbers, and let $\mu,\mu',\mu''$ be
integer-valued sequences. If $\mu'$ agrees with $\mu$ at every index satisfying $K_1$,
$\mu''$ agrees with $\mu'$ at every index satisfying $K_2$, and every index satisfying
$K$ satisfies both $K_1$ and $K_2$, then $\mu''$ agrees with $\mu$ at every index
satisfying $K$.
\end{theorem}

\begin{theorem}[Writing outside the constrained indices]
\lean{Light.SameOn.write}
\label{Light.SameOn.write}
\leanok
\uses{Light.SameOn}
\difficulty{3}
Let $K$ be a predicate on the natural numbers and let $\mu,\mu'$ be integer-valued
sequences. If they agree wherever $K$ holds, and $K$ does not hold at an index $b$, then
for every integer $x$, replacing the value of $\mu'$ at $b$ by $x$ leaves it agreeing
with $\mu$ wherever $K$ holds.
\end{theorem}

\begin{definition}[Writing a finite block into a sequence]
\lean{Light.wrote}
\label{Light.wrote}
\leanok
Given an integer-valued sequence $\mu$, a destination $dst$, a sequence $f$, and a
length $j$, define the resulting sequence to equal $f(a-dst)$ at indices $a$ with
$dst\le a<dst+j$, and to equal $\mu(a)$ at all other indices.
\end{definition}

\begin{theorem}[Writing an Empty Interval]
\lean{Light.wrote_zero}
\label{Light.wrote_zero}
\leanok
\uses{Light.wrote}
\difficulty{3}
Let $\mu,f:\mathbb{N}\to\mathbb{Z}$ and $d\in\mathbb{N}$. Writing zero values of $f$
into $\mu$ starting at $d$ leaves $\mu$ unchanged.
\end{theorem}

\begin{theorem}[Values at Written Indices]
\lean{Light.wrote_done}
\label{Light.wrote_done}
\leanok
\uses{Light.wrote}
\difficulty{3}
Let $\mu,f:\mathbb{N}\to\mathbb{Z}$ and $d,j,i\in\mathbb{N}$, with $i<j$. After writing
the first $j$ values of $f$ into $\mu$ starting at $d$, the value at index $d+i$ is
$f(i)$.
\end{theorem}

\begin{theorem}[Values Outside a Written Interval]
\lean{Light.wrote_rest}
\label{Light.wrote_rest}
\leanok
\uses{Light.Outside, Light.wrote}
\difficulty{3}
Let $\mu,f:\mathbb{N}\to\mathbb{Z}$ and $d,j,a\in\mathbb{N}$. If $a<d$ or $d+j\le a$,
then after writing the first $j$ values of $f$ into $\mu$ starting at $d$, the value at
$a$ remains $\mu(a)$.
\end{theorem}

\begin{theorem}[Extending a Written Interval]
\lean{Light.wrote_succ}
\label{Light.wrote_succ}
\leanok
\uses{Light.wrote, Light.wrote_done}
\difficulty{4}
Let $\mu,f:\mathbb{N}\to\mathbb{Z}$ and $d,j\in\mathbb{N}$. Updating the sequence
obtained by writing the first $j$ values of $f$ into $\mu$ starting at $d$, at index
$d+j$ to the value $f(j)$, gives exactly the sequence obtained by writing the first
$j+1$ values of $f$ starting at $d$.
\end{theorem}

\begin{theorem}[Updating Outside a Written Interval]
\lean{Light.update_wrote}
\label{Light.update_wrote}
\leanok
\uses{Light.Outside, Light.wrote, Light.wrote_rest}
\difficulty{4}
Let $\mu,f:\mathbb{N}\to\mathbb{Z}$, and let $d,j,y\in\mathbb{N}$ and $w\in\mathbb{Z}$.
Suppose that $y<d$ or $d+j\le y$. Updating at $y$ the sequence obtained by writing the
first $j$ values of $f$ into $\mu$ starting at $d$ gives the same sequence as first
updating $\mu$ at $y$ to $w$ and then making that write.
\end{theorem}

\begin{theorem}[A bounded write preserves outside values]
\lean{Light.sameOutside_wrote}
\label{Light.sameOutside_wrote}
\leanok
\uses{Light.SameOutside, Light.wrote, Light.wrote_rest}
\difficulty{1}
Let $\mu$ be a memory function, and let $dst,j,n$ be natural numbers and $f$ an
integer-valued function. If $j\le n$, then writing the values $f(0),\ldots,f(j-1)$ into
the $j$ cells starting at $dst$ leaves memory unchanged at every address outside the
interval from $dst$ through $dst+n-1$.
\end{theorem}
